\documentclass[11pt]{article}

% ---------- layout ----------
\usepackage[margin=1in]{geometry}
\usepackage{microtype}
\usepackage{parskip}
\usepackage{enumitem}
\usepackage{booktabs}
\usepackage{tabularx}
\usepackage{xcolor}
\usepackage[hidelinks]{hyperref}
\usepackage{titlesec}
\usepackage{fancyhdr}

% ---------- style ----------
\definecolor{PolyUBlue}{HTML}{1F4E79}
\definecolor{NoteColor}{HTML}{B54708}
\definecolor{NoteBg}{HTML}{FFF4E5}

\newcommand{\tempnote}[1]{\textcolor{NoteColor}{\textit{#1}}}
\newcommand{\placeholder}[1]{\tempnote{[#1]}}

\pagestyle{fancy}
\fancyhf{}
\lhead{\textcolor{PolyUBlue}{\textbf{DSAI 4207}}}
\rhead{\textcolor{PolyUBlue}{\textbf{Project Proposal}}}
\cfoot{\thepage}
\setlength{\headheight}{14pt}

\titleformat{\section}{\Large\bfseries\color{PolyUBlue}}{\thesection.}{0.5em}{}
\titlespacing*{\section}{0pt}{0.8em}{0.3em}
\titleformat{\subsection}{\normalsize\bfseries\color{PolyUBlue}}{}{0pt}{}
\titlespacing*{\subsection}{0pt}{0.6em}{0.2em}
\setlist[itemize]{leftmargin=1.2em, topsep=0.2em, itemsep=0.15em}
\emergencystretch=1.2em

\begin{document}

{\centering
{\LARGE \textbf{Order vs. Scale: A Controlled Study of SFT/DPO Sequencing\\ Across Data Budgets in Small LM Alignment}}\\[0.2em]
{\large \textcolor{PolyUBlue}{\textit{DSAI 4207 Introduction to Large Language Models --- Course Project Proposal}}}\\[0.2em]
{\normalsize \textit{Concept Note Format}}\\[0.3em]
{\normalsize \textit{2--3 pages (excluding references)}}\\[0.5em]
\par}

\vspace{0.1em}
\textcolor{PolyUBlue}{\rule{\linewidth}{0.8pt}}
\vspace{0.5em}

\section*{Project and Team Details}
\begin{tabularx}{\linewidth}{@{}l X@{}}
\toprule
\textbf{Team members} & PAN Haoren, LIU Shiyu, PAN Yubing --- team of 3 \\
\textbf{Main contact} & LIU Shiyu \\
\textbf{Project type} & (A) Reproduce \& extend (Topic D: SFT and Preference Alignment) \\
\textbf{One-line thesis} & We hypothesize that the standard SFT $\rightarrow$ DPO order yields a higher length-controlled win rate than the reverse order at low data budgets (500 prompts), with this order advantage diminishing as the budget scales to 10{,}000. \\
\bottomrule
\end{tabularx}

\section{Problem Statement and Objective}

Almost every open post-training recipe follows the same pipeline: supervised fine-tuning (SFT) first, then preference optimization such as DPO. The order is inherited from InstructGPT and is rarely questioned, even though the two stages update the model with fundamentally different signals --- imitating one reference answer versus pushing apart a chosen/rejected pair. Meanwhile, annotation budget, not GPU time, is the binding constraint for small teams. Two questions are currently answered only by convention. First, does the \emph{order} of SFT and DPO itself change the outcome when data, compute, and optimization steps are strictly matched? Second, does that answer depend on \emph{how much} data is available?

Recent evidence shows the two stages genuinely interact: under fixed budgets, pure SFT wins below roughly 1k examples, while pure DPO from a base model suffers a cold start (Raghavendra et al., ACL 2025). But only the canonical direction was studied; the reverse order DPO $\rightarrow$ SFT has never been benchmarked across data scales, although theory predicts that a later SFT stage can overwrite preference signals installed earlier. The order effect, if real, should be budget-dependent --- largest exactly in the low-data regime where course-scale and small-team projects operate.

\textbf{Research question.} On a fixed summarization task with a fixed starting checkpoint, how does SFT $\rightarrow$ DPO compare with DPO $\rightarrow$ SFT in both length-controlled win rate and capability tax, and how does this order-attributable gap $\mathrm{OA}(n)$ evolve as the prompt budget $n$ grows from 500 to 10{,}000?

\textbf{Objective.} Measure $\mathrm{OA}(n)$ at $n\in\{500, 2000, 10000\}$ on Reddit TL;DR with Qwen2.5-0.5B under strictly matched data and compute, together with the capability tax each recipe pays.

\textbf{Falsifiable hypothesis (H1).} At $n=500$, SFT $\rightarrow$ DPO achieves a higher length-controlled win rate against the shared starting checkpoint than DPO $\rightarrow$ SFT, with the same direction across three seeds and a bootstrap 95\% CI excluding zero. If $\mathrm{OA}(500)\le 0$, the hypothesis is rejected.

\emph{Secondary hypotheses:} \textbf{H2} predicts $\mathrm{OA}(2000)>0$ but smaller than $\mathrm{OA}(500)$; \textbf{H3} predicts $\mathrm{OA}(10{,}000)\rightarrow 0$ (single seed, descriptive); \textbf{H4} predicts alignment gains are primarily stylistic, with smaller arm gaps on IFEval than on judge LCWR.

\section{Related Work and Baselines}

Our work sits at the intersection of three lines. Stiennon et al.\ established human-preference summarization (TL;DR) as the canonical alignment benchmark; Rafailov et al.'s DPO made preference optimization practical without reinforcement learning. Ouyang et al.\ named the \emph{alignment tax}, which we adopt as a first-class metric. LIMA showed alignment needs surprisingly little data; ORPO argued the two stages should be merged, implying order should not matter. The closest work, \emph{Balancing the Budget} (ACL 2025), studied SFT/DPO allocation under fixed budgets but tested only the canonical direction and task performance alone.

We therefore design four baselines that make our contribution precise: (i)~\textbf{pure SFT} and (ii)~\textbf{pure DPO} from the shared checkpoint, reproducing the low-data findings of Raghavendra et al.; (iii)~the \textbf{canonical SFT $\rightarrow$ DPO} recipe as the reference arm; (iv)~a \textbf{cold-start probe} (DPO from the raw base model) verifying our pipeline reproduces their central phenomenon. All arms share the same base checkpoint, identical prompts, and equal epochs, so comparisons are fair. What is new is the missing cell in the literature: the reverse order benchmarked across budgets, plus the capability-tax dimension in the low-data regime.

\section{Proposed Method}

All main arms start from one shared checkpoint $M_0$: Qwen2.5-0.5B-Base warmed up with 5k human-written general instructions (\texttt{no\_robots}), screened to remove any summarization-like samples so that the 500-prompt budget is genuinely the model's first contact with the task. The budget unit is the number of unique prompts $N$; each arm consumes exactly $N$ prompts, differing only in \emph{how} they are used.

\begin{center}
\begin{tabularx}{\linewidth}{@{}l l X@{}}
\toprule
\textbf{Arm} & \textbf{Recipe} & \textbf{Data usage (budget $N$)} \\
\midrule
G2 & SFT only & $N$ chosen demonstrations \\
G3 & DPO only (ref $=M_0$) & $N$ (chosen, rejected) pairs \\
G4 & SFT $\rightarrow$ DPO & Part A chosen ($0.5N$), then Part B pairs ($0.5N$) \\
G5 & DPO $\rightarrow$ SFT & Part B pairs, then Part A chosen \\
G3$'$ & DPO from raw base & 500 pairs, 1 seed (cold-start probe) \\
\bottomrule
\end{tabularx}
\end{center}

The key control is \emph{mirror binding}: SFT text is taken 100\% from the chosen side of the preference pairs, and G4/G5 see word-identical Part A / Part B data in swapped stage order --- the only difference is sequence. Budgets are \emph{nested subsets} ($500\subset2{,}000\subset10{,}000$), so changing budget never changes which items are seen; Part A/B splits are stratified by length, and all pools are deduplicated against the 200-prompt held-out test set (exact match + MinHash). Two-stage arms merge LoRA weights into the base and fully reset optimizer state between stages. Pre-registered hyperparameters: LoRA $r{=}32,\alpha{=}64$; SFT lr $10^{-4}$; DPO lr $5{\times}10^{-6}$, $\beta{=}0.1$; exactly one epoch per stage, since data efficiency is defined as unique samples seen. Seeds: 42/1337/2024 (10k tier: one reference seed, descriptive only).

\textbf{Decisive experiment.} G4 vs.\ G5 at $n=500$ and $n=2{,}000$, three seeds each: word-identical data, equal steps, only order differs. \emph{Success:} $\mathrm{OA}(500)>0$, consistent across seeds with bootstrap CI excluding zero --- order matters and is budget-dependent, quantitatively confirming sequential-erasure theory; a vanishing $\mathrm{OA}(10{,}000)$ would bound where the convention stops mattering. \emph{Failure:} no significant order gap --- order is then a superstition at small scale, supporting ORPO-style merged objectives; this negative result is equally reportable, and we would analyze which factor (model scale, warmup, budget granularity) masks the predicted effect.

\section{Data, Models, Setup, and Evaluation}

\textbf{Data.} OpenAI \texttt{summarize\_from\_feedback} (Reddit TL;DR): SFT demonstrations and comparison pairs from the same prompt pool; 200 held-out prompts for evaluation; decontamination of all training pools and of the warmup set against test prompts and benchmarks. \textbf{Model.} Qwen2.5-0.5B-Base (open weights, Apache 2.0) with LoRA; decoding is greedy with max 128 new tokens. \textbf{Metrics} (pre-registered): (i)~\emph{gain} --- length-controlled win rate (LCWR) vs.\ $M_0$ judged by GPT-4o-mini with A/B position swapping and a length covariate; (ii)~\emph{cost} --- capability tax on MMLU (500 items) and ARC-Challenge (300 items) scored by log-likelihood to avoid parsing collapse of a 0.5B model; (iii)~\emph{authenticity} --- IFEval with the official programmatic checker, separating real instruction-following gains from stylistic inflation. \textbf{Randomness:} three seeds at 500/2k; conclusions require consistent direction, 1{,}000-sample bootstrap CIs, and McNemar tests on judge pairs. \textbf{Budget:} 29 fine-tuning runs totaling $\sim$10--14 GPU-hours on a single RTX 4090/A10G; judge API cost $\sim\$5$ (3{,}200 short calls). \textbf{Risk and fallback:} if DPO diverges, lower lr to $10^{-6}$ and raise $\beta$ to 0.3; if the 10k tier overruns, it is dropped to descriptive status and inference rests on the fully replicated 500/2k tiers.

\section{Timeline and Roles}
\begin{tabularx}{\linewidth}{@{}l X l@{}}
\toprule
\textbf{Week} & \textbf{Milestone} & \textbf{Owner(s)} \\
\midrule
Week 8--9 & Warmup screening + $M_0$ trained; nested splits, mirror binding, decontamination; end-to-end MVP pipeline & \placeholder{Name A}, \placeholder{Name B} \\
Week 10--11 & All 500/2k runs (24 fine-tunes); 10k reference + cold-start probe; IFEval and capability-tax scoring & \placeholder{Name B}, \placeholder{Name C} \\
Week 12 & Hypothesis tests H1--H4, gain--tax scatter analysis, failure-case analysis; report drafting & \placeholder{Name C}, all \\
Week 13 & Final report/video and presentation rehearsal & All \\
\bottomrule
\end{tabularx}

\vspace{0.4em}
\placeholder{Name A} owns data engineering (warmup screening, nested splits, mirror binding, MinHash decontamination, IFEval harness); \placeholder{Name B} owns the training pipeline ($M_0$ warmup, TRL recipes, two-stage LoRA merge and optimizer reset, scheduling of all 29 runs); \placeholder{Name C} owns evaluation and statistics (judge scripts with position swapping, likelihood scoring, bootstrap/McNemar analysis, figures, report). All members jointly blind-rate 100 judged pairs for human--judge agreement (Cohen's $\kappa$).

\section{Responsible Use and Tool Disclosure}

Qwen2.5 (Apache 2.0), \texttt{no\_robots} (Apache 2.0), and \texttt{summarize\_from\_feedback} (MIT) are used within license; Reddit-derived content is for non-commercial research only; GPT-4o-mini is accessed under OpenAI's API terms. AI assistants (ChatGPT/Claude/Copilot) support brainstorming, debugging, and writing per course policy. All experiments and analyses are our own work.

\section*{References}
\begin{enumerate}[leftmargin=1.4em, itemsep=0.1em]
\item Raghavendra, Kang, Ritter. \emph{Balancing the Budget: Understanding Trade-offs Between Supervised and Preference-Based Finetuning}. ACL 2025. arXiv:2502.11284.
\item Rafailov et al. \emph{Direct Preference Optimization: Your Language Model is Secretly a Reward Model}. NeurIPS 2023.
\item Stiennon et al. \emph{Learning to Summarize from Human Feedback}. NeurIPS 2020.
\item Ouyang et al. \emph{Training Language Models to Follow Instructions with Human Feedback}. NeurIPS 2022.
\item Hong, Lee, Thorne. \emph{ORPO: Monolithic Preference Optimization without Reference Model}. EMNLP 2024.
\item Zhou et al. \emph{Instruction-Following Evaluation for Large Language Models}. arXiv:2311.07911.
\item Yang et al. \emph{Qwen2.5 Technical Report}. arXiv:2412.15115.
\item Hugging Face H4. \emph{No Robots: 10{,}000 Human-Written Instructions for SFT}. \url{https://huggingface.co/datasets/HuggingFaceH4/no_robots}, 2023.
\end{enumerate}

\end{document}
